% !TEX root = ../eigenfaces_algorithm.tex
%=====================================================
\section{交互演示与作图注意事项}

演示逻辑：预处理只做一次——算出 $\bar{\mathbf{a}}$、全部 279 个特征对
（前 50 个来自幂迭代 + 收缩，其余来自 \texttt{eigh}）、缓存全部坐标与各 $k$ 的准确率。
滑块只切 $k$，实时重算 $\mathbf{c}=\mat{V}_k\mT\mathbf{x}$、$\hat{\mathbf{a}}=\mat{V}_k\mathbf{c}+\bar{\mathbf{a}}$，
同屏显示：原图、重建图、前 16 张特征脸，并读出 $R(k)$ 与该 $k$ 的测试准确率。

\textbf{特征脸的显示规范（重要）：}特征向量的分量是实数、有正有负，
必须使用统一色标（灰度中点 $=0$，或按全体分量的全局分位裁剪）。
若每张图各自做 min--max 归一化，16 张图会全都呈现“差不多亮”的假象——
这是初学者报告里最常见的视觉误导。
